\chapter{相关工作}

针对拜占庭容错多智能体协作的研究背景，前期的探索主要集中在以下三个方面。通过分析这些工作的模型设定与处理机制，可以明确在去中心化优化的背景下引入局部代价函数所面临的种种挑战。

\section{基于约简机制的鲁棒迭代一致性}

在对抗环境下的分布式协作中，确保正常节点状态的稳定收敛是首要任务。Vaidya等人在其 2012 年的技术报告中，针对点对点通信模型下的任意有向图，深入探讨了迭代式近似共识的理论边界\upcite{vaidya2012iterative}。该工作通过引入修剪算子——即在节点接收到的状态序列中，排序并严格剔除最大与最小的 $f$ 个值（其中常数 $f$ 代表系统预设的全网拜占庭节点数量上限）——有效地屏蔽了恶意节点通过非一致性攻击注入的伪造数据。该研究从数学上确立了约简图（即从原通信图中移除所有可能的故障节点和与之相连的边以及此后正常节点至多 $f$ 条入边后残存的子图）必须满足包含非空源组件的充分必要条件，为后续处理状态偏差提供了拓扑分析的框架。

与此同时，LeBlanc等人在早期工作中提出了一类基于子序列约简的鲁棒迭代算法\upcite{leblanc2012consensus}。该算法的核心机理在于要求各节点在聚合信息前，主动忽略那些显著偏离自身当前状态的极端观测值，从而保障系统的渐近一致性。为了对网络的抗干扰能力进行量化，该研究定义了基于拓扑凝聚力的 $(r, s)$-鲁棒性指标。在数学上，该指标要求网络中任意两个不相交且集合大小不超过 $s$（设定为局部最多受破坏的节点数）的节点集合，至少有一个集合包含这样的节点：其在自身集合之外拥有至少 $r$ 个邻居（即确保正常节点保留足够多的可靠信息源）。然而，此类工作由于仅针对静态的标量或向量状态进行分析，并未涉及局部输入为成本函数时的梯度更新逻辑，无法直接处理分布式的连续寻优任务。

\section{广播模型下的分布式函数计算}

针对更为复杂的计算任务，部分研究聚焦于在对抗网络中执行特定的函数映射。Sundaram和Hadjicostis提出了一种基于线性迭代策略的分布式函数计算方案，旨在使各节点协作计算出如算术平均值或多项式函数等特定的全局目标\upcite{sundaram2011distributed}。该工作的一个关键假设是采用了广播通信模型，即物理层确保每个节点发送的信息对于其所有入度邻居是完全一致的。这种物理层约束极大限制了拜占庭节点的对抗空间，使其无法针对不同邻居实施具有欺骗性的不一致信息注入。

进一步的理论推导表明，在广播模型假定下，只要网络的顶点连通度(记为 $\kappa$)满足 $\kappa \ge 2f+1$（其中 $f$ 为网络中存在的恶意节点总数），系统即可通过维护状态的历史序列实现对全局函数的精确重构。这里的 $\kappa \ge 2f+1$ 意味着任意两个正常节点之间至少存在 $2f+1$ 条节点不相交的独立路径，从而利用多数表决机制滤除最多 $f$ 条路径上的污染。然而，这种方案的有效性高度依赖于底层协议提供的通信一致性，在通用的点对点网络环境中，广播特性往往难以获得保障。此外，如何在不具备物理层过滤的前提下处理优化过程中由异构局部函数产生的梯度偏差，仍是当前分布式寻优领域亟待攻克的理论难题。

\section{局部有界故障模型下的协作分析}

在探讨故障规模对系统性能的影响时，研究者提出了不同于全局数量限制的故障建模方法。Zhang和Sundaram探讨了信息扩散算法在局部有界故障模型下的鲁棒性\upcite{zhang2011robustness}。在数学表达上，若记正常节点 $i$ 的邻居集合为 $\mathcal{N}_i$，全网拜占庭节点集合为 $\mathcal{A}$，该模型严格约束 $|\mathcal{N}_i \cap \mathcal{A}| \le f$。即假设在任何正常节点的局部邻域内，遭受拜占庭故障的节点数量均不超过预设常数 $f$，而非像传统模型那样仅对全网故障总数施加全局约束（$|\mathcal{A}| \le b$）。这种建模方式能够更准确地刻画超大规模网络中攻击行为的物理分布特性，并据此推导出了实现信息扩散的边界拓扑条件。

尽管局部故障模型简化了大规模拓扑下攻击的复杂度分析，但在实际的分布式优化应用中，算法通常仍需在更具一般性的全局故障模型下保持稳健。此外，现有关于局部故障模型的研究多侧重于信息的扩散速率与拓扑覆盖范围，尚未充分探讨在数据非独立同分布条件下，如何通过严格的数学约束（例如为正常节点分配非平凡的凸系数权重）来防止全局决策被极端观测值或伪装的对抗数据所劫持。因此，在构建具备结构性保护机制的容错优化框架方面，仍存在显著的理论探索空间。
